\documentclass[10pt,a4paper]{amsart}
\usepackage{amsmath,amssymb,amsthm,graphicx}
% Retained minimal preamble for the Stacks Project chapter flat.tex (More on Flatness).
% Verbatim excerpt of the chapter preamble: Xy-pic setup (native preamble lines 14-18) and the native
% theorem declarations with the mathematical macros (native preamble lines 154-217). Copied mechanically;
% nothing added, removed or redefined. The wrapper class and the amsmath/amssymb/amsthm/graphicx base
% packages are supplied by the assembler, so the class/externaldocument part of the native preamble is omitted.
% For commutative diagrams we use Xy-pic
\usepackage[all]{xy}

% We use 2cell for 2-commutative diagrams.
\xyoption{2cell}
% Theorem environments.
%
\theoremstyle{plain}
\newtheorem{theorem}[subsection]{Theorem}
\newtheorem{proposition}[subsection]{Proposition}
\newtheorem{lemma}[subsection]{Lemma}

\theoremstyle{definition}
\newtheorem{definition}[subsection]{Definition}
\newtheorem{example}[subsection]{Example}
\newtheorem{exercise}[subsection]{Exercise}
\newtheorem{situation}[subsection]{Situation}

\theoremstyle{remark}
\newtheorem{remark}[subsection]{Remark}
\newtheorem{remarks}[subsection]{Remarks}

\numberwithin{equation}{subsection}

% Macros
%
\def\lim{\mathop{\mathrm{lim}}\nolimits}
\def\colim{\mathop{\mathrm{colim}}\nolimits}
\def\Spec{\mathop{\mathrm{Spec}}}
\def\Hom{\mathop{\mathrm{Hom}}\nolimits}
\def\Ext{\mathop{\mathrm{Ext}}\nolimits}
\def\SheafHom{\mathop{\mathcal{H}\!\mathit{om}}\nolimits}
\def\SheafExt{\mathop{\mathcal{E}\!\mathit{xt}}\nolimits}
\def\Sch{\mathit{Sch}}
\def\Mor{\mathop{\mathrm{Mor}}\nolimits}
\def\Ob{\mathop{\mathrm{Ob}}\nolimits}
\def\Sh{\mathop{\mathit{Sh}}\nolimits}
\def\NL{\mathop{N\!L}\nolimits}
\def\CH{\mathop{\mathrm{CH}}\nolimits}
\def\proetale{{pro\text{-}\acute{e}tale}}
\def\etale{{\acute{e}tale}}
\def\QCoh{\mathit{QCoh}}
\def\Ker{\mathop{\mathrm{Ker}}}
\def\Im{\mathop{\mathrm{Im}}}
\def\Coker{\mathop{\mathrm{Coker}}}
\def\Coim{\mathop{\mathrm{Coim}}}

% Boxtimes
%
\DeclareMathSymbol{\boxtimes}{\mathbin}{AMSa}{"02}

%
% Macros for moduli stacks/spaces
%
\def\QCohstack{\mathcal{QC}\!\mathit{oh}}
\def\Cohstack{\mathcal{C}\!\mathit{oh}}
\def\Spacesstack{\mathcal{S}\!\mathit{paces}}
\def\Quotfunctor{\mathrm{Quot}}
\def\Hilbfunctor{\mathrm{Hilb}}
\def\Curvesstack{\mathcal{C}\!\mathit{urves}}
\def\Polarizedstack{\mathcal{P}\!\mathit{olarized}}
\def\Complexesstack{\mathcal{C}\!\mathit{omplexes}}
% \Pic is the operator that assigns to X its picard group, usage \Pic(X)
% \Picardstack_{X/B} denotes the Picard stack of X over B
% \Picardfunctor_{X/B} denotes the Picard functor of X over B
\def\Pic{\mathop{\mathrm{Pic}}\nolimits}
\def\Picardstack{\mathcal{P}\!\mathit{ic}}
\def\Picardfunctor{\mathrm{Pic}}
\def\Deformationcategory{\mathcal{D}\!\mathit{ef}}
\usepackage{cleveref}
\crefname{originallemma}{original lemma}{original lemmas}
\Crefname{originallemma}{Original lemma}{Original lemmas}
\makeatletter
\newlabel{divisors-lemma-dominate-admissible-blowups}{{080N}{source}{}{sourcepointer-1}{}}
\newlabel{divisors-lemma-dominate-admissible-blowups@cref}{{[originallemma][0][]080N}{[1][0][]source}{}{}{}}
\makeatother
\crefname{originalremark}{original remark}{original remarks}
\Crefname{originalremark}{Original remark}{Original remarks}
\makeatletter
\newlabel{remark-successive-blowups}{{080Y}{source}{}{sourcepointer-2}{}}
\newlabel{remark-successive-blowups@cref}{{[originalremark][0][]080Y}{[1][0][]source}{}{}{}}
\makeatother
\crefname{originallemma}{original lemma}{original lemmas}
\Crefname{originallemma}{Original lemma}{Original lemmas}
\makeatletter
\newlabel{morphisms-lemma-quasi-affine-finite-type-over-S}{{04II}{source}{}{sourcepointer-3}{}}
\newlabel{morphisms-lemma-quasi-affine-finite-type-over-S@cref}{{[originallemma][0][]04II}{[1][0][]source}{}{}{}}
\makeatother
\crefname{originaltheorem}{original theorem}{original theorems}
\Crefname{originaltheorem}{Original theorem}{Original theorems}
\makeatletter
\newlabel{theorem-flatten-module}{{0815}{source}{}{sourcepointer-4}{}}
\newlabel{theorem-flatten-module@cref}{{[originaltheorem][0][]0815}{[1][0][]source}{}{}{}}
\makeatother
\crefname{originallemma}{original lemma}{original lemmas}
\Crefname{originallemma}{Original lemma}{Original lemmas}
\makeatletter
\newlabel{divisors-lemma-strict-transform-affine}{{080G}{source}{}{sourcepointer-5}{}}
\newlabel{divisors-lemma-strict-transform-affine@cref}{{[originallemma][0][]080G}{[1][0][]source}{}{}{}}
\makeatother
\title{A U-admissible blowup makes the strict transform flat and of finite presentation, for an over-S scheme flat over a quasi-compact open U and of finite type over S}
\author{The Stacks Project Authors}
\date{stacks-project:flat.tex:local-snapshot-20261011; GNU FDL 1.2 or later}
\begin{document}
\maketitle
\noindent\textit{Editorial note.} All delivered mathematical text is the author's own, copied from the pinned Stacks Project chapter file flat.tex; no statement or proof step was rewritten, reordered or removed.\par\medskip
\noindent\textit{Editorial note.} The preamble is the chapter's own retained minimal preamble (sources/stacks-flat/retained-preamble.tex), a verbatim excerpt of the native Xy-pic setup and of the native theorem declarations and mathematical macros; no macro was invented or redefined, and the wrapper class and the amsmath/amssymb/amsthm/graphicx base packages are supplied by the assembler.\par\medskip
\noindent\textit{Editorial note.} The engine of the proof is the flattening theorem of the same chapter, which is retained as a tagged navigation pointer rather than as a counted item; the reduction from the strict transform of the pushforward of the structure sheaf to the strict transform of X is written out in the counted span.\par\medskip
\noindent\textit{Editorial note.} Canonical location of the counted proposition: Stacks tag 081R (sources/stacks-flat/pinned-original-source.tex line 9123).\par\medskip
\noindent\textit{Allowed standard prerequisites (editorial).}\begin{itemize}
\item The flattening theorem for a quasi-coherent module on a scheme of finite type over a quasi-compact quasi-separated base. Same chapter, Theorem labelled theorem-flatten-module (source line 9017), invoked at proof line 9158. Supplies the U-admissible blowup S' -> S making the strict transform of i\_*O\_X flat and of finite presentation; it is the engine of the proof, and the reduction that turns the conclusion about i\_*O\_X into the conclusion about X' is carried out in the counted span.
\item Admissible blowups: dominating a finite family, and behaviour of strict transforms. Divisors, Lemma labelled divisors-lemma-dominate-admissible-blowups (Stacks tag 080N) with Remark labelled remark-successive-blowups (source line 8611), and Divisors, Lemma labelled divisors-lemma-strict-transform-affine (Stacks tag 080G), invoked at proof lines 9148, 9166. Supply the blowup dominating all the local solutions and the commutation of strict transform with pushforward along an affine morphism.
\item An affine scheme of finite type over S embeds into an affine space over S. Morphisms, Lemma labelled morphisms-lemma-quasi-affine-finite-type-over-S (Stacks tag 04II), invoked at proof line 9154. Supplies the immersion j : X -> A\textasciicircum{}n\_S and the quasi-compact open V in which j stays a closed immersion.
\end{itemize}\medskip
\ifdefined\hypertarget\hypertarget{sourcepointer-1}{}\fi
\noindent\textit{Source reference.} Original lemma 080N: sources/stacks-chow/tags.txt, source lines 10188--10188. Cross-chapter standard prerequisite located in the project tag list staged in the local snapshot; invoked at proof line 9148 of the counted proof. It is retained only as a navigation pointer; every load-bearing step is inside the delivered span.\par\medskip
\ifdefined\hypertarget\hypertarget{sourcepointer-2}{}\fi
\noindent\textit{Source reference.} Original remark 080Y: sources/stacks-flat/pinned-original-source.tex, source lines 8611--8611. Same-chapter pointer into the author environment remark at source line 8611; invoked at proof line 9149 of the counted proof. It is retained only as a navigation pointer; every load-bearing step is inside the delivered span.\par\medskip
\ifdefined\hypertarget\hypertarget{sourcepointer-3}{}\fi
\noindent\textit{Source reference.} Original lemma 04II: sources/stacks-chow/tags.txt, source lines 5904--5904. Cross-chapter standard prerequisite located in the project tag list staged in the local snapshot; invoked at proof line 9154 of the counted proof. It is retained only as a navigation pointer; every load-bearing step is inside the delivered span.\par\medskip
\ifdefined\hypertarget\hypertarget{sourcepointer-4}{}\fi
\noindent\textit{Source reference.} Original theorem 0815: sources/stacks-flat/pinned-original-source.tex, source lines 9017--9017. Same-chapter pointer into the author environment theorem at source line 9017; invoked at proof line 9158 of the counted proof. It is retained only as a navigation pointer; every load-bearing step is inside the delivered span.\par\medskip
\ifdefined\hypertarget\hypertarget{sourcepointer-5}{}\fi
\noindent\textit{Source reference.} Original lemma 080G: sources/stacks-chow/tags.txt, source lines 10181--10181. Cross-chapter standard prerequisite located in the project tag list staged in the local snapshot; invoked at proof line 9166 of the counted proof. It is retained only as a navigation pointer; every load-bearing step is inside the delivered span.\par\medskip
\begin{lemma}
\label{lemma-flat-after-blowing-up}
Let $S$ be a quasi-compact and quasi-separated scheme.
Let $X$ be a scheme over $S$.
Let $U \subset S$ be a quasi-compact open.
Assume
\begin{enumerate}
\item $X \to S$ is of finite type and quasi-separated, and
\item $X_U \to U$ is flat and locally of finite presentation.
\end{enumerate}
Then there exists a $U$-admissible blowup $S' \to S$ such that
the strict transform of $X$ is flat and of finite presentation
over $S'$.
\end{lemma}

\begin{proof}
Since $X \to S$ is quasi-compact and quasi-separated by assumption,
the strict transform of $X$ with respect to a blowing up $S' \to S$
is also quasi-compact and quasi-separated. Hence to prove the lemma
it suffices to find a $U$-admissible blowup such that the strict
transform is flat and locally of finite presentation.
Let $X = W_1 \cup \ldots \cup W_n$ be a finite affine open covering.
If we can find a $U$-admissible blowup $S_i \to S$ such that the
strict transform of $W_i$ is flat and locally of finite presentation,
then there exists a $U$-admissible blowing up $S' \to S$ dominating
all $S_i \to S$ which does the job (see
Divisors, Lemma \ref{divisors-lemma-dominate-admissible-blowups};
see also Remark \ref{remark-successive-blowups}).
Hence we may assume $X$ is affine.

\medskip\noindent
Assume $X$ is affine. By
Morphisms, Lemma \ref{morphisms-lemma-quasi-affine-finite-type-over-S}
we can choose an immersion $j : X \to \mathbf{A}^n_S$ over $S$.
Let $V \subset \mathbf{A}^n_S$ be a quasi-compact open subscheme
such that $j$ induces a closed immersion $i : X \to V$ over $S$. Apply
Theorem \ref{theorem-flatten-module}
to $V \to S$ and the quasi-coherent module $i_*\mathcal{O}_X$
to obtain a $U$-admissible blowup $S' \to S$ such that the strict
transform of $i_*\mathcal{O}_X$ is flat over $S'$ and of finite presentation
over $\mathcal{O}_{V \times_S S'}$. Let $X'$ be the strict transform
of $X$ with respect to $S' \to S$. Let $i' : X' \to V \times_S S'$
be the induced morphism.
Since taking strict transform commutes with pushforward along affine
morphisms (Divisors, Lemma \ref{divisors-lemma-strict-transform-affine}),
we see that $i'_*\mathcal{O}_{X'}$ is flat over $S$ and of
finite presentation as a $\mathcal{O}_{V \times_S S'}$-module.
This implies the lemma.
\end{proof}

\end{document}
